% Title page of the PDF guides.
% Included by `make wiki-pdf` through pandoc -H; filled in and called by the guide's own file
% title-<guide>.tex (pandoc -B).
%
% Why not through the pandoc variables title/subtitle/date: the input is the wiki's sidebar,
% and a YAML header in a wiki page would show on GitHub. Version, date and commit come from
% the wiki checkout instead (the Makefile fills them in).

\usepackage{array}

\newcommand{\DocTitle}{}
\newcommand{\DocSubtitle}{}
\newcommand{\DocLines}{}
\newcommand{\DocFootnote}{}

\newcommand{\doctitle}[1]{\renewcommand{\DocTitle}{#1}}
\newcommand{\docsubtitle}[1]{\renewcommand{\DocSubtitle}{#1}}
\newcommand{\docfootnote}[1]{\renewcommand{\DocFootnote}{#1}}
% One row of the table of particulars. Can be called several times; the order of the calls
% is the order on the page.
\newcommand{\docline}[2]{\appto\DocLines{#1 & #2\\}}

\definecolor{titlerule}{HTML}{1F3864}
\definecolor{titlegrey}{HTML}{4A4A4A}

\newcommand{\TitlePage}{%
  \begin{titlepage}
  \thispagestyle{empty}
  \sffamily
  \noindent
  {\color{titlegrey}\footnotesize\MakeUppercase{fsbpyr\quad\textbar\quad github.com/devitops-com/fsbpyr}}\par
  \vspace{2pt}
  % \hrule instead of \rule: a rule inside a paragraph runs over the text block by
  % \parindent (Overfull \hbox). \hrule is set in vertical mode and is exactly \hsize wide.
  {\color{titlerule}\hrule height 1.2pt}

  \vspace{5.5cm}

  \noindent{\fontsize{30}{34}\selectfont\bfseries\DocTitle\par}

  \vspace{10pt}
  \noindent{\color{titlegrey}\fontsize{13}{18}\selectfont\raggedright\DocSubtitle\par}

  \vspace{1.4cm}
  {\color{titlerule}\hrule height 0.6pt}
  \vspace{10pt}

  \noindent\footnotesize
  \begin{tabular}{@{}>{\bfseries}p{4.2cm}p{\dimexpr\linewidth-5cm\relax}@{}}
    \DocLines
  \end{tabular}

  \vfill

  {\color{titlerule}\hrule height 0.6pt}
  \vspace{6pt}
  \noindent{\color{titlegrey}\scriptsize\DocFootnote\par}
  \end{titlepage}
  \clearpage
}
